% ============================================================================
% Fichier  : partie4/P04_C13_modeles_formels.tex
% Titre    : Modèles Formels de Sécurité
% Auteur   : MINKA MI NGUIDJOÏ Thierry Emmanuel
% Version  : 1.0 -- novembre 2026
% ============================================================================
% Base : brouillon « Livre IV » (Bell-LaPadula, Biba, Chinese Wall) et
% « Livre III » (modèles de contrôle d'accès). Sources : Lampson (1971) ;
% Bell et LaPadula (1976) ; Biba (1977) ; Clark et Wilson (1987) ; Brewer et
% Nash (1989) ; Sandhu et al. (1996) ; NIST SP 800-162 (2014).
% ============================================================================

\chapter{Modèles Formels de Sécurité}
\label{chap:modeles}
\chaptermark{Modèles formels}

\epigraph{%
    « Une politique que l'on ne sait pas énoncer\\
    est une politique que l'on ne sait pas vérifier. »%
}{Principe du cours}

\niveauChapitre{Intermédiaire}{Chapitres~\ref{chap:proprietes}
et~\ref{chap:actifs}~; notions de logique.}

\begin{boxobjectifs}
\begin{itemize}
    \item Représenter des droits par une matrice d'accès, des listes et
          des capacités.
    \item Énoncer les règles de Bell-LaPadula et de Biba, et ce qu'elles
          protègent.
    \item Expliquer les principes de Clark-Wilson et la muraille de Chine.
    \item Distinguer les politiques DAC, MAC, RBAC et ABAC.
    \item Choisir un modèle selon la propriété à protéger.
\end{itemize}
\end{boxobjectifs}

\bigskip

% ============================================================================
\section*{Question directrice : quelle règle aurait empêché un compte de tout lire~?}
% ============================================================================

À Garoua, l'attaquant, devenu administrateur du domaine, a pu lire tous
les dossiers et effacer les sauvegardes. Le système appliquait une règle
simple~: l'administrateur peut tout. Les modèles formels montrent qu'on
peut énoncer d'autres règles, plus fines, et les vérifier.

\medskip
\textbf{Plan du chapitre.}
\begin{enumerate}
    \item La matrice d'accès (section~\ref{sec:c13-matrice}).
    \item Bell-LaPadula~: la confidentialité (section~\ref{sec:c13-blp}).
    \item Biba~: l'intégrité (section~\ref{sec:c13-biba}).
    \item Clark-Wilson et la muraille de Chine (section~\ref{sec:c13-commerce}).
    \item DAC, MAC, RBAC, ABAC (section~\ref{sec:c13-politiques}).
\end{enumerate}

\separateur

% ============================================================================
\section{La Matrice d'Accès}
\label{sec:c13-matrice}
% ============================================================================

B.~Lampson a proposé de représenter la protection par une matrice dont
les lignes sont les \emph{sujets} (utilisateurs, processus), les colonnes
les \emph{objets} (fichiers, tables), et les cases les droits~\cite{lampson1971}.

\begin{tableaucours}{L{3.2cm}L{2.8cm}L{2.8cm}L{3.2cm}}{%
    \tableauHeader{Sujet} &
    \tableauHeader{Factures} &
    \tableauHeader{Paiements} &
    \tableauHeader{Sauvegardes}
}
Mme~OWONA (saisie) & lire & lire, écrire & aucun \\
Validatrice & lire & lire, valider & aucun \\
Service de sauvegarde & lire & lire & écrire \\
Administrateur & lire, écrire & lire, écrire & lire, écrire, effacer \\
\end{tableaucours}
\vspace{-0.8em}
\begin{center}{\small\itshape Tableau~13.1. Matrice d'accès (extrait) de LiZbiZ}\end{center}

La matrice se stocke par colonnes (\emph{listes de contrôle d'accès},
attachées aux objets) ou par lignes (\emph{capacités}, détenues par les
sujets). Les systèmes de fichiers courants utilisent des listes de
contrôle d'accès.

% ============================================================================
\section{Bell-LaPadula : la Confidentialité}
\label{sec:c13-blp}
% ============================================================================

Conçu pour les systèmes militaires, le modèle de D.~Bell et L.~LaPadula
attribue à chaque sujet une \emph{habilitation} et à chaque objet une
\emph{classification}, ordonnées (non classifié < confidentiel < secret)
\cite{bell1976}.

\begin{boxdef}
\termecle{Bell-LaPadula}~: deux règles. \emph{Pas de lecture vers le
haut}~: un sujet ne lit pas un objet de niveau supérieur à son
habilitation. \emph{Pas d'écriture vers le bas} (propriété~\(\star\))~:
un sujet n'écrit pas dans un objet de niveau inférieur, pour qu'une
information secrète ne fuie pas vers un document moins protégé.
\end{boxdef}

Le modèle protège la confidentialité, et seulement elle~: rien n'empêche
un sujet de bas niveau d'écrire n'importe quoi vers le haut.

% ============================================================================
\section{Biba : l'Intégrité}
\label{sec:c13-biba}
% ============================================================================

K.~Biba a proposé le modèle dual, pour l'intégrité~\cite{biba1977}.

\begin{boxdef}
\termecle{Biba}~: \emph{pas de lecture vers le bas}~: un sujet ne lit pas
une information de moindre intégrité, qui pourrait le corrompre~; \emph{pas
d'écriture vers le haut}~: un sujet n'écrit pas dans un objet d'intégrité
supérieure.
\end{boxdef}

\begin{boxscenario}{Opération GAROUA. Les sauvegardes selon Biba}
\textbf{Analyse} (confiance élevée). Les sauvegardes sont l'objet de plus
haute intégrité~: c'est d'elles que l'on repart. Selon Biba, un sujet de
moindre intégrité (un compte d'administration du domaine, exposé à
Internet par l'accès distant) ne devrait pas pouvoir y écrire, ni
l'effacer. Un système de sauvegarde qui n'accepte que des ajouts
(stockage immuable), administré par des comptes distincts, applique cette
règle.
\end{boxscenario}

% ============================================================================
\section{Clark-Wilson et la Muraille de Chine}
\label{sec:c13-commerce}
% ============================================================================

\subsection{Clark-Wilson}

D.~Clark et D.~Wilson ont montré que les entreprises protègent surtout
l'intégrité, par deux principes~: les \emph{transactions bien formées}
(on ne modifie une donnée que par un programme qui la garde cohérente) et
la \emph{séparation des tâches} (aucune personne ne réalise seule une
opération sensible)~\cite{clark1987}. Les données contrôlées ne se
modifient que par des procédures certifiées, et des procédures de
vérification contrôlent régulièrement leur cohérence.

\subsection{La muraille de Chine}

D.~Brewer et M.~Nash ont formalisé la politique des cabinets de conseil~:
un consultant qui a eu accès aux données d'une entreprise ne peut plus
accéder à celles de ses concurrentes~\cite{brewer1989}. Les droits
changent selon l'historique des accès~: c'est une politique contre les
conflits d'intérêts.

% ============================================================================
\section{DAC, MAC, RBAC, ABAC}
\label{sec:c13-politiques}
% ============================================================================

\begin{tableaucours}{L{2.2cm}L{4.8cm}L{5.2cm}}{%
    \tableauHeader{Politique} &
    \tableauHeader{Qui décide des droits} &
    \tableauHeader{Exemple}
}
DAC & Le propriétaire de l'objet, à sa discrétion & Partage d'un dossier
Windows \\
MAC & Le système, selon des étiquettes imposées & Bell-LaPadula~; SELinux \\
RBAC & L'organisation, par rôle métier~\cite{sandhu1996} & Le rôle
« comptable fournisseurs » donne les droits nécessaires \\
ABAC & Des règles sur les attributs du sujet, de l'objet et du
contexte~\cite{nist800162} & « lecture autorisée si service = finance et
heure ouvrée et poste géré » \\
\end{tableaucours}
\vspace{-0.8em}
\begin{center}{\small\itshape Tableau~13.2. Quatre familles de politiques
de contrôle d'accès}\end{center}

\begin{boxalert}
\textbf{Violation typique.} Créer un rôle par personne.

\smallskip
\textbf{Règle~:} un rôle correspond à une fonction métier, pas à un
individu. Une organisation de 150~personnes doit pouvoir se décrire avec
quelques dizaines de rôles, revus chaque année~; sinon, le RBAC dégénère
en DAC déguisé.
\end{boxalert}

\begin{boxobjet}{M-13 · Séparation des droits d'administration}
\textbf{Objectif~:} qu'aucun compte ne puisse à la fois administrer la
production et effacer les sauvegardes. \textbf{Modèle~:} Biba et
séparation des tâches de Clark-Wilson. \textbf{Mise en œuvre~:} comptes
d'administration distincts pour le domaine et pour les sauvegardes~;
stockage de sauvegarde immuable. \textbf{Vérification~:} tentative
d'effacement d'une sauvegarde avec un compte d'administration du domaine,
qui doit échouer.
\end{boxobjet}

\begin{boxmetier}{Développeur}
\textbf{Implanter un RBAC.} Le développeur vérifie le rôle à chaque
requête (médiation complète), refuse par défaut, et journalise les
décisions de refus.
\end{boxmetier}

\begin{boxmetier}{Auditeur de sécurité}
\textbf{Tester la séparation des tâches.} L'auditeur croise la matrice
des droits avec les paires de droits incompatibles (saisir et valider un
paiement), et liste les comptes qui cumulent.
\end{boxmetier}

\separateur

\begin{boxresume}
\textbf{Ce qu'il faut retenir de ce chapitre~:}
\begin{itemize}
    \item \textbf{La matrice d'accès} relie sujets, objets et droits~; on
          la stocke en listes de contrôle ou en capacités.
    \item \textbf{Bell-LaPadula} protège la confidentialité (pas de lecture
          vers le haut, pas d'écriture vers le bas).
    \item \textbf{Biba} protège l'intégrité (pas de lecture vers le bas,
          pas d'écriture vers le haut).
    \item \textbf{Clark-Wilson}~: transactions bien formées et séparation
          des tâches~; \textbf{Brewer-Nash}~: pas d'accès à deux
          concurrents.
    \item \textbf{DAC, MAC, RBAC, ABAC} diffèrent par qui décide des droits.
\end{itemize}
\end{boxresume}

\begin{boxinfo}
\textbf{Points de contrôle de l'élève.} Je sais~: construire une matrice
d'accès~; énoncer les règles de Bell-LaPadula et de Biba~; expliquer les
deux principes de Clark-Wilson~; distinguer DAC, MAC, RBAC et ABAC~;
choisir un modèle selon la propriété à protéger.
\end{boxinfo}

\bigskip

\begin{boxexo}[title={\ding{45}\quad Exercice~13.1. Bell-LaPadula \niveaubase}]
Un sujet habilité « confidentiel » veut lire un document « secret », puis
écrire dans un document « non classifié ». Que dit le modèle pour chaque
opération~?
\end{boxexo}

\begin{boxexo}[title={\ding{45}\quad Exercice~13.2. Biba \niveaubase}]
Pourquoi Biba interdit-il de lire vers le bas~? Donnez un exemple tiré
d'une entreprise.
\end{boxexo}

\begin{boxexo}[title={\ding{45}\quad Exercice~13.3. Quelle politique~? \niveaubase}]
Classez en DAC, MAC, RBAC ou ABAC~: un utilisateur partage un fichier avec
un collègue~; le système refuse l'accès hors du réseau de l'entreprise~;
le rôle « caissier » donne l'accès à la caisse~; des étiquettes de
classification imposées par le système.
\end{boxexo}

\begin{boxexo}[title={\ding{45}\quad Exercice~13.4. La matrice \niveaumoyen}]
Construisez la matrice d'accès de l'agence de la BANQUE SANAGA~: guichetier,
chef d'agence, auditeur interne~; objets~: comptes clients, journal des
opérations, paramètres de plafonds.
\end{boxexo}

\begin{boxexo}[title={\ding{45}\quad Exercice~13.5. Listes ou capacités \niveaumoyen}]
Révoquer tous les droits d'un salarié qui part est-il plus simple avec
des listes de contrôle d'accès ou avec des capacités~? Et retirer l'accès
de tous à un fichier~?
\end{boxexo}

\begin{boxexo}[title={\ding{45}\quad Exercice~13.6. Clark-Wilson \niveaumoyen}]
Décrivez le paiement d'un fournisseur chez LiZbiZ comme une transaction
bien formée, avec sa séparation des tâches et sa procédure de vérification.
\end{boxexo}

\begin{boxexo}[title={\ding{45}\quad Exercice~13.7. Muraille de Chine \niveaumoyen}]
AUDIX CONSEIL audite LiZbiZ et une autre entreprise de logistique.
Comment appliquer la muraille de Chine à ses auditeurs~?
\end{boxexo}

\begin{boxexo}[title={\ding{45}\quad Exercice~13.8. Les deux propriétés \niveaumoyen}]
Peut-on appliquer en même temps Bell-LaPadula et Biba avec les mêmes
niveaux~? Qu'en résulte-t-il pour la circulation de l'information~?
\end{boxexo}

\begin{boxexo}[title={\ding{45}\quad Exercice~13.9. Une règle ABAC \niveauexpert}]
Écrivez trois règles ABAC qui auraient limité l'Opération GAROUA (attributs
du sujet, de la ressource, du contexte).
\end{boxexo}

\begin{boxexo}[title={\ding{45}\quad Exercice~13.10. Des rôles pour LiZbiZ \niveauexpert}]
Proposez une douzaine de rôles RBAC pour LiZbiZ, avec leurs droits
principaux et les paires de rôles incompatibles.
\end{boxexo}

\bigskip

\begin{boxinfo}
\textbf{Pour aller plus loin.}
\begin{itemize}
    \item D.~Bell et L.~LaPadula~\cite{bell1976}~; K.~Biba~\cite{biba1977}.
    \item D.~Clark et D.~Wilson~\cite{clark1987}~; D.~Brewer et
          M.~Nash~\cite{brewer1989}.
    \item NIST SP~800-162 (ABAC)~\cite{nist800162}.
\end{itemize}
\end{boxinfo}

\begin{boxinfo}
\textbf{Transition.} Les modèles disent quelles règles appliquer~;
l'architecture dit où les placer. C'est l'objet du
Chapitre~\ref{chap:architecture}.
\end{boxinfo}
